\medskip
\noindent\emph{Check of the supplied three-puncture calculation.}
We have checked LM3b~H.6--H.8 in its supplied
\texttt{sections/Higher.md}. For H.7, let
$G=F(a,b)$, $I_1=\langle a\rangle$,
$I_2=\langle b\rangle$, and
$I_3=\langle(ab)^{-1}\rangle$.
For each finite-index normal $N$, its strict quotient category has
open group $Q_N=G/N$ and crossing stabilizers $I_{i,N}$, the images
of $I_i$. Its nerve is the homotopy pushout obtained by attaching
three cones on $BI_{i,N}$ to $BQ_N$, by the same crossing-bar
decomposition used in B.3. For a finite constant abelian group $A$,
the filtered colimit over $N$ of these cochain triangles computes
continuous cohomology of $\widehat G$ and the three completed
$I_i$: every continuous cochain factors through a finite quotient,
and filtered colimits of abelian groups are exact. Each $I_i$
has its full induced profinite topology, since its generator is
primitive in $F(a,b)$ and admits a retraction to $\mathbb Z$.
The free and procyclic degree-one resolutions identify these
continuous cohomology groups with the discrete ones. The two
remaining maps are exactly
\[
 \begin{aligned}
 A\oplus A^3&\longrightarrow A^3,&
 (c,d_1,d_2,d_3)&\longmapsto(c-d_1,c-d_2,c-d_3),\\
 A^2&\longrightarrow A^3,&(s,t)&\longmapsto(s,t,-s-t).
 \end{aligned}
\]
The first has diagonal kernel $A$ and is surjective. The second
is injective with cokernel $A$, identified by the sum of the three
coordinates. Hence the filtered nerve cohomology is $A$ in degrees
zero and two and zero elsewhere; the latter cokernel is precisely
the oriented peripheral relation and hence the sphere class.
This verifies H.7 in its finite-coefficient, filtered-colimit form,
without a stagewise cohomology-vanishing assertion. In particular
no one finite stage works uniformly for all finite coefficients:
for a prime $\ell\nmid|Q_N|$, its positive-degree
$\mathbb F_\ell$-cohomology vanishes, including degree two.
H.8's upgrade from these constant coefficients to all finite
homotopy values is separately checked in E.12 and F.43.

\medskip
\noindent\emph{The bounded contractible-object comparison.}
The supplied CATHOT~D.1 represents $[NC,Y]$, for categories $C$
of a specified size bound, by a category of bounded contractible
simplicial sets over $Y$, with all maps over $Y$ retained. That
bounded resolution is a useful comparison of constructions, but
the pieces in B.1 are not individually contractible. Already a
geometric punctured disc has group $\mathbb Z$, or
$\widehat{\mathbb Z}(1)$ for finite coverings, and its classifying
type has nonzero degree-one cohomology. The mixed topos instead
retains the underived crossing biset, whose effective stabilizer is
$K_s$. B.3 proves exactly what is lost on replacing the full local
term $BH_s$ by $BK_s$. Thus the bounded representability result
does not imply the mixed-topos comparison and cannot remove its
local-inertia condition.
